\documentclass[
  aps,
  prd,
  reprint,
  superscriptaddress,
  floatfix,
  nofootinbib,
  longbibliography
]{revtex4-2}

\usepackage[T1]{fontenc}
\usepackage[utf8]{inputenc}
\usepackage{amsmath,amssymb,amsfonts,amsthm,mathtools}
\usepackage{microtype}
\theoremstyle{plain}
\newtheorem{theorem}{Theorem}
\newtheorem{proposition}[theorem]{Proposition}
\newtheorem{corollary}[theorem]{Corollary}
\newtheorem{criterion}[theorem]{Criterion}
\newtheorem{law}[theorem]{Law}
\theoremstyle{definition}
\newtheorem{definition}[theorem]{Definition}
\theoremstyle{remark}
\newtheorem{remark}[theorem]{Remark}

\usepackage{hyperref}

\hypersetup{
  colorlinks=true,
  linkcolor=blue,
  citecolor=blue,
  urlcolor=blue
}

\newcommand{\dd}{\mathrm d}
\newcommand{\D}{\mathcal D}
\newcommand{\vis}{\mathrm{vis}}
\newcommand{\wind}{\mathrm w}
\newcommand{\eff}{\mathrm{eff}}
\newcommand{\crit}{\mathrm{crit}}
\newcommand{\GI}{\mathrm{GI}}
\newcommand{\Ok}{\Omega_K}
\newcommand{\Ow}{\Omega_{\mathrm w}}
\newcommand{\Oq}{\Omega_q}
\newcommand{\Aq}{\mathcal A_q}
\newcommand{\Aw}{\mathcal A_{\mathrm w}}
\newcommand{\Rcl}{\mathcal R}
\DeclareMathOperator{\Vol}{Vol}
\DeclareMathOperator{\Spec}{Spec}
\DeclareMathOperator{\Tr}{Tr}
\newcommand{\id}{\mathbf 1}
\newcommand{\Z}{\mathbb Z}
\newcommand{\R}{\mathbb R}
\newcommand{\eps}{\epsilon}
\newcommand{\qhat}{\widehat q}
\newcommand{\NF}{N_F}
\newcommand{\eF}{e_F}
\newcommand{\PG}{P_\Gamma}
\newcommand{\Hn}{\mathcal H_n}
\newcommand{\LCDM}{\Lambda\mathrm{CDM}}

\begin{document}
\raggedbottom

\title{\texorpdfstring{Einstein-Residual Closure Standard:\
Invariant Dark-Sector Audits, Normal Forms, and Topology-Likelihood Boundaries}
{Einstein-Residual Closure Standard: Invariant Dark-Sector Audits, Normal Forms, and Topology-Likelihood Boundaries}}

\author{Andrei~Preece}
\email{andre.preece@web.de}
\affiliation{Independent Researcher, Berlin, Germany}

\author{Boris~Batenin}
\email{batenin@web.de}
\affiliation{Independent Researcher, Berlin, Germany}

\date{May 16, 2026 -- v28 universal-standard synthesis}

\begin{abstract}
A dark-sector model should not be evaluated before its geometric residual,
visible-sector convention, homogeneous normal form, microscopic realization,
amplitude selector, perturbative transfer map, and likelihood benchmark have
been separately declared.  We formulate this requirement as an
Einstein-residual closure standard: a convention-explicit audit protocol rather
than a new local law of nature.  For a fixed spacetime metric, a declared
Einstein-side convention, and a specified visible-sector subtraction, the
Einstein residual
\[
T^\D_{\mu\nu}=\frac{c^4}{8\pi G}G_{\mu\nu}-T^{\vis}_{\mu\nu}
\]
is a covariantly fixed reconstruction target rather than a source assumption.
In the homogeneous locally-FLRW sector the residual admits the lossless atlas
\[
\Aq=-p_\D,
\qquad
\Aw=a^3(\eps_\D+p_\D),
\]
and is exactly minimal dust plus fixed vacuum if and only if
\[
\frac{\dd\Aq}{\dd\ln a}=0,
\qquad
\frac{\dd\Aw}{\dd\ln a}=0,
\qquad
\Aq,\Aw\ge0 .
\]
This theorem is algebraic, not an abundance-origin mechanism.  Its minimal
background realization is a conserved Brown current with linear energy law plus
a fixed four-form vacuum branch.  Brown-current dynamics closes the dust
equation of state but not the dust abundance; the four-form closes \(w=-1\) but
not the vacuum amplitude.  Compact causal geometry supplies conformal null-reach
invariants, and compact topology enters through quotient projectors, spectra,
observable covariance, and observer priors; neither selects \(\Omega_c\),
\(\Omega_\Lambda\), or \(\Lambda_{\eff}\) by itself.  We make the standard
operational by introducing an admissibility vector, hierarchy no-go rules,
closure-defect and likelihood ledgers, and a minimal benchmark interface.  The
resulting principle is: no dark-sector claim without residual reconstruction,
closure-defect test, microscopic ledger, amplitude ledger, transfer ledger, and
matched benchmark likelihood.
\end{abstract}

\maketitle

\section{Claim boundary}

This paper gives a conditional homogeneous stress-tensor classification.  The
classification is algebraic rather than dynamical: it does not derive the
observed dark-sector abundances, does not select a flux branch, and does not
claim detection of compact topology.  No abundance-origin claim is made.  Once
the metric \(g_{\mu\nu}\), the Einstein-side convention, and the visible-sector
subtraction \(T^{\vis}_{\mu\nu}\) are fixed, the Einstein residual is fixed.  For
that same convention and visible subtraction, if this residual is homogeneous
and locally FLRW, then it has the minimal dust-like plus fixed-vacuum
stress-tensor form if and only if
\begin{equation}
-p_\D=\text{constant}\ge0,
\qquad
a^3(\eps_\D+p_\D)=\text{constant}\ge0 .
\label{eq:claim_boundary_closure}
\end{equation}
Microscopic current and four-form realizations may impose further amplitude
selection rules; the closure theorem itself does not supply them.  The
corresponding stress tensor is locally indistinguishable from cold dark matter
plus \(\Lambda\) with the same amplitudes and curvature.  Any distinction must
therefore come from compact spectral projection, quotient covariance,
perturbative null conditions, or an additional amplitude-selection rule.

The paper uses a ledger hierarchy.  Level I is residual reconstruction from
\((g_{\mu\nu},T^{\vis}_{\mu\nu})\).  Level II is the homogeneous normal form in
\((\Aq,\Aw)\).  Level III is Brown-current realization of the dust equation of
state.  Level IV is fixed four-form realization of the vacuum equation of
state.  Level V is compact causal reach, which constrains null accessibility but
adds no new homogeneous source.  Level VI is perturbative source behavior, such
as dust sound speed and anisotropic-stress null conditions.  Level VII is a
topology-likelihood claim, where the quotient projector has been propagated into
a masked, noisy observable covariance.  This manuscript supplies the Level-I/II
normal-form theorem and disciplined completion ledgers for Levels III--VII; it
does not perform a completed topology likelihood or amplitude selection.

The local residual theorem applies to any FLRW curvature class once the
homogeneous metric, the visible subtraction, and the Einstein-side convention
are declared.  The compact topology ledger used below is narrower: it is
specialized to compact orientable spherical space forms \(S_A^3/\Gamma\), except
where Euclidean compact-topology papers are cited explicitly as covariance-level
methodological analogues.

The paper is complete only at this declared level: homogeneous residual
closure, compact-background bookkeeping, perturbative null conditions, and
amplitude-selection boundaries.  It is not a compact Boltzmann solver, a
likelihood analysis, or a microscopic dark-sector completion.


\section{Einstein-residual closure standard}

The preceding boundary can be elevated from a single model statement to a
reproducible standard.  The standard is not a new local law of nature.  It is a
claim-control protocol for dark-sector models: each model must identify which
part of the argument is reconstruction, which part is microscopic realization,
which part is amplitude selection, and which part is observational likelihood.

\begin{law}[Primitive residual law]
Given \((g_{\mu\nu},T^{\vis}_{\mu\nu})\) and an Einstein-side convention, the
only primitive dark-sector reconstruction target is
\begin{equation}
T^\D_{\mu\nu}[g,V]
=
\frac{c^4}{8\pi G}G_{\mu\nu}[g]-T^{\vis}_{\mu\nu}[V].
\label{eq:standard_primitive_residual}
\end{equation}
It is a residual tensor, not automatically a particle species, vacuum branch,
topological sector, or modified-gravity mechanism.
\end{law}

\begin{law}[Subtraction covariance]
If the declared visible subtraction is changed by a homogeneous perfect-fluid
piece \(X_{\mu\nu}=\eps_X U_\mu U_\nu/c^2+p_Xh_{\mu\nu}\), then
\begin{equation}
T^\D_{\mu\nu}\mapsto T^\D_{\mu\nu}-X_{\mu\nu},
\label{eq:standard_subtraction_covariance_tensor}
\end{equation}
and the homogeneous atlas transforms affinely as
\begin{equation}
\Aq\mapsto \Aq+p_X,
\qquad
\Aw\mapsto \Aw-a^3(\eps_X+p_X).
\label{eq:standard_subtraction_covariance_atlas}
\end{equation}
Consequently, no residual-closure claim is invariant until the visible
subtraction and Einstein-side convention have been declared.
\end{law}

\begin{law}[FLRW atlas completeness]
In the homogeneous locally-FLRW perfect-fluid class,
\begin{equation}
(\eps_\D,p_\D)\longleftrightarrow(\Aq,\Aw)
\end{equation}
with inverse
\begin{equation}
p_\D=-\Aq,
\qquad
\eps_\D=a^{-3}\Aw+\Aq .
\label{eq:standard_atlas_inverse}
\end{equation}
Thus any further distinction between two models with the same atlas is not a
homogeneous-background distinction; it must arise from microscopic realization,
amplitude selection, perturbations, topology, or likelihood.
\end{law}

\begin{law}[Minimal closure law]
The homogeneous residual is exactly minimal dust plus fixed vacuum on a connected
scale-factor interval \(I\) if and only if
\begin{equation}
\frac{\dd\Aq}{\dd\ln a}=0,
\qquad
\frac{\dd\Aw}{\dd\ln a}=0,
\qquad
\Aq,\Aw\ge0 .
\label{eq:standard_minimal_closure_law}
\end{equation}
In that case
\begin{equation}
T^\D_{\mu\nu}
=
\eps_{\wind,0}a^{-3}\frac{U_\mu U_\nu}{c^2}
-
\eps_qg_{\mu\nu}.
\label{eq:standard_minimal_tensor}
\end{equation}
This is an iff statement for the homogeneous stress tensor; it is not an
amplitude-origin theorem.
\end{law}

\begin{law}[Closure-defect invariant]
For any reference density \(\eps_*>0\), define
\begin{equation}
\mathfrak C_I[\D;\eps_*]
:=
\sup_{a\in I}
\frac{
\left(\frac{\dd\Aq}{\dd\ln a}\right)^2+
\left(a^{-3}\frac{\dd\Aw}{\dd\ln a}\right)^2
}{\eps_*^2} .
\label{eq:standard_closure_defect}
\end{equation}
Then \(\mathfrak C_I=0\) if and only if the adopted residual is exactly minimal
dust plus fixed vacuum on \(I\).  The value depends on the chosen norm and
reference density, but the zero set does not.
\end{law}

\begin{law}[Brown-current dust law]
For a Brown-current realization with \(N^\mu=nu^\mu\), \(\nabla_\mu N^\mu=0\),
and energy law \(F(n)\),
\begin{equation}
\eps_J=F(n),
\qquad
p_J=nF'(n)-F(n).
\end{equation}
The pressureless condition is equivalent to
\begin{equation}
nF'(n)=F(n),
\end{equation}
whose regular nonzero local solution is
\begin{equation}
F(n)=\mu_{\wind}c^2n.
\label{eq:standard_brown_dust_law}
\end{equation}
Thus dust behavior is linear current energy, not topology.
\end{law}

\begin{law}[Current-amplitude no-go]
A topological degree \(Q\in\Z\) labels a sector.  It becomes a Brown-current
number only after the additional branch rule
\begin{equation}
N_J=|Q|
\label{eq:standard_branch_rule_NJQ}
\end{equation}
has been imposed.  Even then,
\begin{equation}
\Ow=
\frac{\mu_{\wind}c^2|Q|}{\Vol(\Sigma_{A_0})\eps_{\crit,0}}
\label{eq:standard_current_amplitude_nogo}
\end{equation}
requires the non-topological scale \(\mu_{\wind}\).  Therefore \((Q,\Gamma)\)
by themselves do not determine \(\Ow\).
\end{law}

\begin{law}[Four-form vacuum law]
A four-form branch in four spacetime dimensions obeys, on a fixed connected
branch,
\begin{equation}
F_{\mu\nu\rho\sigma}=f\varepsilon_{\mu\nu\rho\sigma},
\qquad
\dd f=0,
\end{equation}
and therefore has vacuum stress tensor
\begin{equation}
T^q_{\mu\nu}=-\eps_q g_{\mu\nu}.
\label{eq:standard_four_form_vacuum}
\end{equation}
It closes the equation of state \(w=-1\), not the observed amplitude
\(\eps_q\).
\end{law}

\begin{law}[Flux-lattice separation]
The continuous fixed-vacuum statement
\begin{equation}
\eps_q=\frac12\qhat^2,
\qquad
\qhat\in\R,
\end{equation}
is distinct from the added flux-lattice rule
\begin{equation}
\qhat=\NF\eF,
\qquad
\NF\in\Z .
\end{equation}
Flux quantization requires a normalization or membrane sector; it is not a
consequence of compact spatial topology alone.
\end{law}

\begin{law}[Topological decoupling]
There is no dimensionally complete topology-only map
\begin{equation}
(\Gamma,Q,\NF)\not\Rightarrow(\Ow,\Oq,\Lambda_{\eff}).
\label{eq:standard_topological_decoupling}
\end{equation}
Any amplitude selector must factor through additional non-topological data,
for example
\begin{equation}
\mathfrak S_{\rm amp}:
(\Gamma,Q,\NF;\mu_{\wind},\eF,A_0,H_0,\mathcal B)
\longrightarrow
(\Ow,\Oq).
\label{eq:standard_amp_selector}
\end{equation}
\end{law}

\begin{law}[Compact spectral projector]
For \(\Sigma_\Gamma=S^3/\Gamma\), physical fields are projected cover fields,
\begin{equation}
\Phi(\gamma x)=\Phi(x),
\qquad
\gamma\in\Gamma,
\end{equation}
with projector
\begin{equation}
\PG=\frac1{|\Gamma|}\sum_{\gamma\in\Gamma}\gamma^*.
\label{eq:standard_projector_law}
\end{equation}
Topology enters through \(\PG\), \(\Spec_\Gamma\), \(m_\Gamma(n)\), observable
covariance, and observer priors, not through a new local homogeneous equation
of state.
\end{law}

\begin{law}[Compact causal accessibility]
For a round simply connected cover,
\begin{equation}
\Delta\chi(t_1,t_2)=\int_{t_1}^{t_2}\frac{c\,\dd t}{a(t)A_0}.
\end{equation}
Full antipodal reach on \(S^3\) requires \(\Delta\chi\ge\pi\).  For a quotient,
the observer-dependent sufficient reach condition is
\begin{equation}
\Delta\chi\ge D_\Gamma(x_o),
\qquad
D_\Gamma(x_o)=\sup_{y\in\Sigma_\Gamma}d_\Gamma(x_o,y).
\label{eq:standard_causal_quotient_law}
\end{equation}
Compact causality constrains null reach; it does not select energy density.
\end{law}

\begin{law}[Perturbative degeneracy]
Within a declared data model, the same background amplitudes, visible sector,
primordial statistics, and compact transfer functions imply the same observable
covariance:
\begin{equation}
\begin{aligned}
&\text{same }(\eps_{\wind,0},\eps_q),\quad
\text{same visible sector},
\\
&\text{same }\Xi^\Gamma,\quad
\text{same }\Delta_{\ell n}^{X,\Gamma}
\Longrightarrow
\text{same }{\rm Cov}^{XY,\Gamma} .
\end{aligned}
\label{eq:standard_perturbative_degeneracy}
\end{equation}
No perturbative observable can distinguish two models after this covariance has
been matched.
\end{law}

\begin{law}[Benchmark likelihood]
A topology claim must compare against a matched benchmark,
\begin{equation}
\Delta\ln\mathcal L
=
\ln\mathcal L_\Gamma-
\ln\mathcal L_{\rm bench},
\label{eq:standard_benchmark_likelihood}
\end{equation}
where the benchmark has the same local expansion history, curvature, visible
sector, amplitudes, mask, noise, and foreground conventions.  Otherwise the
comparison confounds topology with background or data-model choices.
\end{law}

The standard can be summarized as a single admissibility principle:
\begin{quote}
A dark-sector claim is admissible only after its Einstein residual, closure
atlas, microscopic realization, amplitude selector, perturbative transfer map,
and benchmark likelihood have been separately declared and invariantly tested.
\end{quote}

\section{Universal admissibility invariants}

The closure standard becomes useful beyond this manuscript only if it separates
universal bookkeeping from model-specific dynamics.  The universal object is not
``dark matter,'' ``dark energy,'' or ``topology.''  It is the ordered declaration
\begin{equation}
\mathfrak M=
(g_{\mu\nu},\mathcal E,T^{\vis}_{\mu\nu},I,\mathcal S_{\rm micro},
\mathcal S_{\rm amp},\mathcal T_{\rm pert},\mathcal D,\mathcal L_{\rm bench}),
\label{eq:model_declaration_v28}
\end{equation}
where \(\mathcal E\) denotes the Einstein-side convention, \(I\) the
scale-factor interval, \(\mathcal S_{\rm micro}\) the microscopic source ledger,
\(\mathcal S_{\rm amp}\) the amplitude selector, \(\mathcal T_{\rm pert}\) the
perturbative transfer map, \(\mathcal D\) the declared data vector, and
\(\mathcal L_{\rm bench}\) the matched benchmark likelihood.  A model that omits
one entry may still be mathematically meaningful, but it is not entitled to the
corresponding stronger claim.

\begin{definition}[Admissibility vector]
For a declared model \(\mathfrak M\), define the binary or graded vector
\begin{equation}
\mathbf A(\mathfrak M)=
(A_{\rm R},A_{\rm H},A_{\rm J},A_q,A_{\rm amp},A_{\rm T},A_{\rm L}),
\label{eq:admissibility_vector_v28}
\end{equation}
where \(A_{\rm R}\) records residual reconstruction, \(A_{\rm H}\) homogeneous
closure, \(A_{\rm J}\) current or matter-source realization, \(A_q\) vacuum
realization, \(A_{\rm amp}\) amplitude selection, \(A_{\rm T}\) perturbative
transfer closure, and \(A_{\rm L}\) likelihood-grade comparison.  The minimal
binary convention is
\begin{equation}
A_i\in\{0,1\},
\end{equation}
with \(A_i=1\) only when the corresponding row of the audit table has been
explicitly supplied.  A numerical implementation may replace the binary entries
by calibrated defects or Bayes factors, but the order of the entries is fixed.
\end{definition}

\begin{law}[Admissibility monotonicity]
A higher-level claim cannot be inferred from a lower-level entry of
\(\mathbf A\).  In particular,
\begin{align}
A_{\rm H}=1&\not\Rightarrow A_{\rm amp}=1,\nonumber\\
A_{\rm amp}=1&\not\Rightarrow A_{\rm L}=1,\nonumber\\
A_{\rm T}=1&\not\Rightarrow A_{\rm L}=1 .
\label{eq:admissibility_monotonicity_v28}
\end{align}
Thus a background closure theorem, a source EFT, an amplitude selector, and a
likelihood claim are logically distinct achievements.  This is the basic
anti-overclaim invariant of the standard.
\end{law}

\begin{law}[Hierarchy no-go]
Within the declared FLRW residual standard,
\begin{equation}
\begin{array}{rcl}
\text{closure atlas} &\not\Rightarrow& \text{microphysics},\\
\text{Brown current} &\not\Rightarrow& \text{dust abundance},\\
\text{four-form branch} &\not\Rightarrow& \text{vacuum amplitude},\\
\text{compact quotient} &\not\Rightarrow& \text{dark-sector amplitudes},\\
\text{causal reach} &\not\Rightarrow& \text{energy density},\\
\text{projected modes} &\not\Rightarrow& \text{topology detection}.
\end{array}
\label{eq:hierarchy_nogo_v28}
\end{equation}
Each implication becomes meaningful only after the missing ledger is supplied:
source action, mass or flux scale, branch prior, transfer functions, data vector,
and matched likelihood.
\end{law}

\begin{theorem}[Standard completeness at the declared level]
For homogeneous locally-FLRW residuals with fixed
\((g_{\mu\nu},\mathcal E,T^{\vis}_{\mu\nu})\), the pair \((\Aq,\Aw)\), the
closure defect \(\mathfrak C_I\), and the sign conditions \(\Aq,\Aw\ge0\)
exhaust all background information relevant to whether the residual is minimal
dust plus fixed vacuum.  Any further question belongs to exactly one of the
following ledgers: source realization, amplitude selection, perturbative
transfer, compact causal/spectral structure, or likelihood comparison.
\label{thm:standard_completeness_v28}
\end{theorem}

\begin{proof}
The map \((\eps_\D,p_\D)\leftrightarrow(\Aq,\Aw)\) is algebraically invertible,
so there is no remaining homogeneous perfect-fluid information after the atlas
is known.  The zero set of \(\mathfrak C_I\) is exactly the constant-atlas
condition, and the sign conditions are exactly \(\eps_{\wind,0},\eps_q\ge0\).
Brown-current, four-form, topology, perturbation, and likelihood data do not
alter this local algebraic classification; they determine whether the algebraic
form is physically realized, amplitude-selected, observationally transferred,
or statistically favored.  Hence the listed ledgers are complete for the objects
introduced by the standard, while making no claim about theories outside the
standard's declared scope.
\end{proof}

\begin{definition}[Claim classes]
A manuscript-level statement is assigned to one of the following claim classes:
\begin{equation}
\begin{array}{ll}
\mathcal C_{\mathrm R}: & \text{residual reconstruction},\\
\mathcal C_{\mathrm H}: & \text{homogeneous closure},\\
\mathcal C_{\mathrm S}: & \text{source realization},\\
\mathcal C_{\mathrm A}: & \text{amplitude selection},\\
\mathcal C_{\mathrm P}: & \text{perturbative transfer},\\
\mathcal C_{\mathrm G}: & \text{global causal/spectral topology},\\
\mathcal C_{\mathrm L}: & \text{likelihood-grade evidence}.
\end{array}
\label{eq:claim_classes_v28}
\end{equation}
The standard forbids promoting a statement from \(\mathcal C_i\) to a stronger
class unless the additional invariant objects of that stronger class have been
specified.
\end{definition}

\begin{criterion}[Gold-standard success criterion]
A dark-sector framework reaches gold-standard status only when it supplies a
reproducible tuple
\begin{equation}
\begin{aligned}
\mathcal G_{\mathrm{std}}=(
&\mathrm{theorem},\ \mathrm{audit\ table},\ \mathrm{code\ interface},\\
&\mathrm{public\ examples},\ \mathrm{benchmark\ likelihood}) .
\end{aligned}
\label{eq:gold_standard_tuple_v28}
\end{equation}
The theorem prevents conceptual overclaiming, the table makes the claim class
visible, the code interface makes closure defects reproducible, public examples
show how standard models pass or fail, and the benchmark likelihood prevents
confusing topology, amplitudes, and data-model choices.
\end{criterion}



\section{Placement in the current topology program}

The present result is deliberately narrower than a topology search.  Current
cosmic-topology work has moved from matched-circle exclusions toward harmonic
covariance, quotient eigenmodes, observer-position dependence, tensor-mode
covariance, and robustness against non-standard primordial spectra.  Planck
found no evidence for the specific compact topologies and anisotropic geometries
tested in its matched-circle and Bayesian pipelines, principally flat toroidal
or slab topologies and Bianchi VII\(_h\) geometries; this is not an exclusion of
all compact spherical quotients~\cite{PlanckTopology}.  Spherical lens-space
circle constraints must be read with their own \(p,q\), curvature-scale, and
observer-domain assumptions~\cite{SahaLens}.

The COMPACT covariance program is used here as a methodological analogue unless
a cited paper is explicitly spherical.  Its Euclidean analyses show the required
likelihood-level structure: matched-circle limits do not exhaust topology
constraints, quotient boundary conditions alter eigenmodes and off-diagonal
harmonic covariance, observer position can matter, tensor sectors have their own
covariance, and non-standard primordial spectra can amplify or suppress
detectability~\cite{COMPACTPartI,COMPACTIIa,COMPACTIIb,COMPACTIIIb,
COMPACTIIc}.  These statements support the ledger structure below; they are not
used as a completed likelihood pipeline for every \(S^3/\Gamma\).

Quantum-field effects on compact spaces define a separate stress-energy source.
A topology-induced Casimir stress is not an amplitude selector for the minimal
homogeneous branch.  If such a stress produces anisotropic backreaction, it
leaves the homogeneous locally-FLRW residual class rather than closing it from
inside~\cite{NegroCasimir}.

This paper therefore fixes the boundary of a background theorem.  Its invariant
input is only
\[
(g_{\mu\nu},\;T^{\vis}_{\mu\nu},\;\text{Einstein-side convention},\;\Gamma),
\]
and its output is a local stress-tensor normal form together with the conditions
under which a compact quotient can enter only through volume, spectra,
covariance, and observer priors.  It does not replace the covariance program;
it supplies a convention-safe background ledger for it.

\section{Einstein residual and convention covariance}

Let \((M,g_{\mu\nu})\) be a four-dimensional spacetime, and let
\(T^{\vis}_{\mu\nu}\) denote the stress tensor assigned to the declared visible
sector.  Define
\begin{equation}
\boxed{
T^\D_{\mu\nu}:=
\frac{c^4}{8\pi G}G_{\mu\nu}-T^{\vis}_{\mu\nu} .
}
\label{eq:primitive_residual_v5}
\end{equation}
A bare cosmological-constant term is included in \(T^\D_{\mu\nu}\) unless the
Einstein-side convention is explicitly changed to put \(G_{\mu\nu}+\Lambda
 g_{\mu\nu}\) on the left-hand side.

\begin{proposition}[Residual uniqueness and convention covariance]
If
\begin{equation}
G_{\mu\nu}=\frac{8\pi G}{c^4}\left(T^{\vis}_{\mu\nu}+S_{\mu\nu}\right),
\end{equation}
then \(S_{\mu\nu}=T^\D_{\mu\nu}\).  If the visible subtraction is changed by
\(X_{\mu\nu}\), then
\begin{equation}
T^{\vis}_{\mu\nu}\to T^{\vis}_{\mu\nu}+X_{\mu\nu},
\qquad
T^\D_{\mu\nu}\to T^\D_{\mu\nu}-X_{\mu\nu} .
\end{equation}
If the Einstein side is changed from \(G_{\mu\nu}\) to
\(G_{\mu\nu}+\Lambda_0g_{\mu\nu}\) while the visible subtraction is kept fixed,
then
\begin{equation}
T^\D_{\mu\nu}\to T^\D_{\mu\nu}+\frac{c^4\Lambda_0}{8\pi G}g_{\mu\nu} .
\end{equation}
This is a change in the definition of the Einstein-side residual, not the usual
operation of moving a positive cosmological constant from the left-hand side to
the right-hand side, where it is represented by
\(T^\Lambda_{\mu\nu}=-(c^4\Lambda/8\pi G)g_{\mu\nu}\).
\end{proposition}

\begin{proposition}[Visible-subtraction covariance of the atlas]
\label{prop:atlas_subtraction_covariance_v26}
Suppose the homogeneous visible subtraction is changed by a homogeneous
perfect-fluid tensor
\begin{equation}
X_{\mu\nu}=\eps_X\frac{U_\mu U_\nu}{c^2}+p_X h_{\mu\nu},
\end{equation}
so that \(T^{\vis}_{\mu\nu}\to T^{\vis}_{\mu\nu}+X_{\mu\nu}\) and
\(T^\D_{\mu\nu}\to T^\D_{\mu\nu}-X_{\mu\nu}\).  Then
\begin{equation}
\eps_\D\to \eps_\D-\eps_X,
\qquad
p_\D\to p_\D-p_X,
\end{equation}
and the closure-pair coordinates transform affinely as
\begin{equation}
\boxed{
\Aq\to \Aq+p_X,
\qquad
\Aw\to \Aw-a^3(\eps_X+p_X) .
}
\label{eq:atlas_affine_subtraction_v26}
\end{equation}
Consequently, the statement that a residual is dust plus fixed vacuum is not an
absolute statement about the metric alone.  It is an invariant statement only
after the Einstein-side convention and the visible-sector subtraction have been
declared.
\end{proposition}

We use standard differential-geometric conventions~\cite{Wald}.  The contracted Bianchi identity gives
\begin{equation}
\nabla^\mu T^\D_{\mu\nu}=-\nabla^\mu T^{\vis}_{\mu\nu} .
\label{eq:bianchi_residual_v5}
\end{equation}
Thus separate residual conservation is not automatic; it follows only when the
visible sector is separately conserved in the adopted convention.  This point
matters because the minimal dust-plus-vacuum source branch is separately
conserved.

We use a timelike four-velocity \(U^\mu\) normalized by
\(U^\mu U_\mu=-c^2\) and the projector
\begin{equation}
h_{\mu\nu}=g_{\mu\nu}+\frac{U_\mu U_\nu}{c^2} .
\label{eq:projector_v5}
\end{equation}
For a perfect fluid with physical energy density \(\eps\) and pressure \(p\),
\begin{equation}
T_{\mu\nu}=\eps\frac{U_\mu U_\nu}{c^2}+p h_{\mu\nu} .
\label{eq:perfect_fluid_v5}
\end{equation}

\section{Locally-FLRW reconstruction}

On the round cover, use
\begin{equation}
\dd s^2=-c^2\dd t^2+a^2(t)A_0^2\dd\Omega_3^2,
\qquad
\Ok=-\frac{c^2}{H_0^2A_0^2}<0,
\label{eq:closed_metric_v5}
\end{equation}
where \(A_0\) is the present curvature radius of the simply connected cover.
For a quotient \(S^3_{A_0}/\Gamma\), the local metric is unchanged while the
volume and global spectrum are restricted.  Let
\begin{equation}
E(a)=\frac{H(a)}{H_0},
\qquad
\eps_{\crit,0}=\frac{3H_0^2c^2}{8\pi G} .
\end{equation}
The local Friedmann equations reconstruct a homogeneous residual as
\begin{align}
\eps_\D(a)
&=\eps_{\crit,0}\left[E^2(a)-\Ok a^{-2}\right]
-\eps_{\vis}(a),
\label{eq:eps_reconstruction_v5}\\
p_\D(a)
&=-\frac{\eps_{\crit,0}}{3}\biggl[
2E(a)\frac{\dd E}{\dd\ln a}
\nonumber\\
&\hspace{1.0cm}
+3E^2(a)-\Ok a^{-2}
\biggr]-p_{\vis}(a).
\label{eq:p_reconstruction_v5}
\end{align}
These are kinematical reconstructions from the chosen FLRW metric and subtraction convention.  They do not assert residual field equations or any microscopic source model.

\begin{theorem}[Locally-FLRW residual reduction]
After a fixed visible subtraction, every homogeneous and isotropic symmetric
residual tensor is uniquely of the form
\begin{equation}
T^\D_{\mu\nu}=\eps_\D(a)\frac{U_\mu U_\nu}{c^2}+p_\D(a)h_{\mu\nu} .
\label{eq:flrw_residual_tensor_v5}
\end{equation}
If the residual is separately conserved, then
\begin{equation}
\frac{\dd\eps_\D}{\dd\ln a}+3(\eps_\D+p_\D)=0 .
\label{eq:residual_conservation_v5}
\end{equation}
\end{theorem}

\begin{proof}
Local homogeneity and isotropy leave only \(U_\mu U_\nu\) and the spatial
projector \(h_{\mu\nu}\) as invariant symmetric rank-two tensors.  Taking the
covariant divergence in FLRW gives Eq.~\eqref{eq:residual_conservation_v5}.
\end{proof}

\section{Closure-pair atlas and minimal branch}

Define the closure variables
\begin{equation}
\boxed{
\Aq(a):=-p_\D(a),
\qquad
\Aw(a):=a^3[\eps_\D(a)+p_\D(a)] .
}
\label{eq:closure_pair_v5}
\end{equation}
They are a lossless reparametrization of the homogeneous residual:
\begin{equation}
p_\D(a)=-\Aq(a),
\qquad
\eps_\D(a)=a^{-3}\Aw(a)+\Aq(a) .
\label{eq:atlas_inverse_v5}
\end{equation}

\begin{proposition}[Completeness of the homogeneous residual atlas]
\label{prop:atlas_completeness_v26}
Within the homogeneous locally-FLRW perfect-fluid residual class, the pair
\((\Aq,\Aw)\) is a complete set of coordinates.  Any scalar diagnostic depending
only on the homogeneous residual stress tensor and the declared scale factor
factors through \((\Aq,\Aw)\).  Therefore no additional homogeneous scalar
criterion can distinguish a residual from dust plus fixed vacuum once
\(\Aq\) and \(\Aw\) are known; all further distinctions must be microscopic,
perturbative, topological, or observational.
\end{proposition}

If the residual is separately conserved, then
\begin{equation}
\frac{\dd\Aq}{\dd\ln a}+a^{-3}\frac{\dd\Aw}{\dd\ln a}=0 .
\label{eq:atlas_flow_v5}
\end{equation}
Thus the general separately conserved homogeneous residual is a one-function
flow in \((\Aq,\Aw)\) space.  The minimal dust-plus-vacuum branch is the
constant nonnegative submanifold.

\begin{theorem}[Homogeneous dust-plus-vacuum residual classification]
\label{thm:minimal_closure_v5}
Let \(I\subset\R_+\) be a connected scale-factor interval, and let
\((\eps_\D,p_\D)\in C^1(I)\) be the homogeneous locally-FLRW Einstein residual
after a fixed visible-sector subtraction.  Then, on $I$, the homogeneous local
stress tensor is representable as
\begin{equation}
T^\D_{\mu\nu}
=
\eps_{\wind,0}a^{-3}\frac{U_\mu U_\nu}{c^2}
-
\eps_q g_{\mu\nu},
\qquad
\eps_{\wind,0},\eps_q\ge0,
\label{eq:minimal_tensor_v5}
\end{equation}
if and only if \(\Aq\) and \(\Aw\) are constant and nonnegative on \(I\).  In
that case
\begin{equation}
\begin{gathered}
\eps_q=\Aq,\qquad \eps_{\wind,0}=\Aw,\\
\eps_\D(a)=\eps_{\wind,0}a^{-3}+\eps_q,
\qquad
p_\D(a)=-\eps_q .
\end{gathered}
\label{eq:closure_solution_v5}
\end{equation}
\end{theorem}

\begin{proof}
If Eq.~\eqref{eq:minimal_tensor_v5} holds, then
\(-p_\D=\eps_q\) and \(a^3(\eps_\D+p_\D)=\eps_{\wind,0}\).  Conversely, if
\(\Aq\) and \(\Aw\) are constant, then Eq.~\eqref{eq:atlas_inverse_v5} gives
\(p_\D=-\Aq\) and \(\eps_\D=\Aw a^{-3}+\Aq\), which is precisely
Eq.~\eqref{eq:minimal_tensor_v5}.  Nonnegativity is exactly the condition
\(\eps_q,\eps_{\wind,0}\ge0\).
\end{proof}

Although this equivalence is algebraically elementary, its role is diagnostic:
after the metric, Einstein-side convention, and visible subtraction are fixed,
it gives the necessary and sufficient closure tests for interpreting the
homogeneous residual as pressureless matter plus fixed vacuum.  The theorem is
therefore a convention-explicit closure diagnostic rather than a dynamical
derivation.  The ``if and only if'' is an iff statement for the homogeneous
stress tensor.
A microscopic branch realizes the constants only if its charge, mass scale,
flux normalization, branch labels, and quotient volume admit those amplitudes.
In particular, a discrete flux lattice can make some otherwise admissible
constant values of \(\eps_q\) unavailable unless the flux unit or branch
spectrum is correspondingly specified.

The invariant null tests are
\begin{equation}
\begin{gathered}
\Rcl_q:=\frac{\dd\Aq}{\dd\ln a}=0,
\qquad
\Rcl_{\wind}:=\frac{\dd\Aw}{\dd\ln a}=0,\\
\Aq,\Aw\ge0 .
\end{gathered}
\label{eq:null_tests_v5}
\end{equation}
For data or reconstruction pipelines it is useful to compress the same null
condition into a nonnegative defect.  For any fixed reference density
\(\eps_*>0\), define
\begin{equation}
\mathfrak C_I[\D;\eps_*]
:=
\sup_{a\in I}
\frac{
\left(\frac{\dd\Aq}{\dd\ln a}\right)^2+
\left(a^{-3}\frac{\dd\Aw}{\dd\ln a}\right)^2
}{\eps_*^2}.
\label{eq:closure_defect_v20}
\end{equation}
The numerical value depends on the chosen norm and on \(\eps_*\), but its zero
set does not:
\begin{align}
\mathfrak C_I[\D;\eps_*]=0
&\Longleftrightarrow
\Rcl_q=\Rcl_{\wind}=0,
\nonumber\\
&\Longleftrightarrow
T^\D_{\mu\nu}\ \text{has the minimal form on }I .
\label{eq:closure_defect_zero_v20}
\end{align}

Equivalently, a reconstruction pipeline is fully specified by the ordered map
\begin{equation}
\begin{aligned}
(E(a),\eps_{\vis},p_{\vis},\Ok)
&\longrightarrow (\eps_\D,p_\D)\\
&\longrightarrow (\Aq,\Aw)\\
&\longrightarrow (\Rcl_q,\Rcl_{\wind},\mathfrak C_I).
\end{aligned}
\label{eq:closure_pipeline_v26}
\end{equation}
This makes the theorem operational: failures of closure are assigned to a
specific visible subtraction and expansion history rather than to topology by
default.

A failure of these tests falsifies the minimal homogeneous branch under the
adopted visible subtraction; it does not falsify compact topology.  No separate-conservation assumption is required for the algebraic iff statement.
However, when the constants are realized, the resulting dust-plus-vacuum residual
is separately conserved.  Hence, by Eq.~\eqref{eq:bianchi_residual_v5}, a
successful minimal decomposition also implies separate conservation of the
declared visible sector in the same convention.  Failure of visible-sector
separate conservation is therefore already a failure of the adopted minimal
source decomposition, not an independent topology effect.

\begin{corollary}[Local \(\LCDM\) degeneracy]
On any local FLRW chart, the minimal branch is exactly the stress tensor of a
pressureless component plus a cosmological-constant component with the same
amplitudes and curvature.  No homogeneous observable depending only on
\(H(a)\), distances, and the same curvature can distinguish the minimal branch
from CDM plus \(\Lambda\).  A CDM plus \(\Lambda\) model placed on the same
compact quotient, with the same projected perturbation covariance, the same
initial data, and the same amplitudes, remains degenerate with the minimal
branch at this level.
\end{corollary}

\section{Current and four-form representation}

One economical microscopic representation of the minimal stress tensor is a
pressureless conserved-current branch plus a fixed four-form branch.  The current branch may carry a protected
degree label, but topology alone does not give a dust stress tensor.

The degree label is standard in algebraic topology~\cite{Hatcher}, while the
pressureless current representative is a Brown-fluid effective
description~\cite{BrownFluid}.  A topological degree may label a sector, but by
itself it does not imply a homogeneous dust current:
\begin{equation}
Q\in\Z\not\Rightarrow N^\mu=n(a)u^\mu\not\Rightarrow p_{\wind}=0 .
\label{eq:q_not_dust_v5}
\end{equation}
Let \(\Sigma\) be compact and oriented, let \(\phi:\Sigma\to S^3\), and
normalize the target volume form by \(\int_{S^3}\widehat\Omega=1\).  The degree is
\begin{equation}
Q=\int_\Sigma \phi^*\widehat\Omega\in\Z .
\label{eq:degree_v5}
\end{equation}
The dust result requires a homogeneous coarse-grained Brown-current
representative and a linear energy density
\begin{equation}
\eps_{\wind}(n)=\mu_{\wind}c^2n .
\label{eq:linear_current_v5}
\end{equation}
More generally, a barotropic current effective theory with
\(\eps_J=F(n)\) has
\begin{equation}
p_J=nF'(n)-F(n),
\qquad
\frac{v_{s,J}^2}{c^2}:=\frac{\dd p_J}{\dd\eps_J}
=\frac{nF''(n)}{F'(n)} .
\label{eq:current_atlas_v5}
\end{equation}
The minimal pressureless branch is the special case \(F(n)=\mu_{\wind}c^2n\),
after any additive vacuum offset has been removed and assigned to the vacuum
sector.  Then \(p_J=0\) and \(v_{s,J}=0\).

A compact-quotient specialization of the Brown current can be written as
\begin{equation}
S_{\rm J}
=\frac1c\int_{\mathbb R\times\Sigma}
\left[-\sqrt{-g}\,F(n)
+{\mathcal J}^{\mu}
\bigl(\partial_\mu\vartheta+\alpha_A\partial_\mu\beta^A\bigr)
\right]\dd^4x,
\label{eq:brown_action_compact_v13}
\end{equation}
where \({\mathcal J}^{\mu}=\sqrt{-g}\,n u^\mu\) is the densitized current,
\(u^\mu u_\mu=-1\), and in this subsection \(u^\mu:=U^\mu/c\) is the
dimensionless unit four-velocity compatible with \(U^\mu U_\mu=-c^2\).  Moreover,
\begin{equation}
n=\frac{1}{\sqrt{-g}}
\sqrt{-g_{\mu\nu}{\mathcal J}^{\mu}{\mathcal J}^{\nu}} .
\label{eq:brown_n_v13}
\end{equation}
The Clebsch potentials enforce current conservation and advected labels; they do
not add a topological amplitude.  Variation gives the perfect-fluid form
\begin{equation}
\begin{aligned}
T^{\rm J}_{\mu\nu}
&=(\eps_{\rm J}+p_{\rm J})u_\mu u_\nu+p_{\rm J}g_{\mu\nu},\\
\eps_{\rm J}&=F(n),
\qquad
p_{\rm J}=nF'(n)-F(n).
\end{aligned}
\label{eq:brown_stress_v13}
\end{equation}
Thus the current action is specialized to a compact quotient once
\((A_0,\Gamma,Q,F)\) are declared, but its amplitude still depends on the
non-topological energy law \(F\).  The integer \(Q\) labels a topological sector;
it becomes a Brown-current number only after an additional branch rule is
imposed.

\begin{criterion}[Current-EFT stability ledger]
For a barotropic Brown-current realization of the residual, a conservative
stability domain is
\begin{equation}
\begin{gathered}
F(n)\ge0,
\qquad
F'(n)>0,\\
0\le \frac{v_s^2}{c^2}=\frac{\dd p}{\dd\eps}=
\frac{nF''(n)}{F'(n)}\le1 .
\end{gathered}
\label{eq:current_stability_v13}
\end{equation}
The first two conditions give positive energy and positive inertial enthalpy,
\(\eps+p=nF'(n)>0\); the last condition gives nonnegative and subluminal
adiabatic response.  Pure pressureless current behavior requires the special stable line
\begin{equation}
F(n)=\mu_{\wind}c^2n,
\qquad
p_{\wind}=0,
\qquad
v_{s,\wind}^2=0,
\label{eq:linear_current_line_v13}
\end{equation}
after any additive vacuum offset has been removed and assigned consistently to
the vacuum sector.  More generally, \(F(n)=\mu_{\wind}c^2n+C\) gives \(p_J=-C\),
namely dust plus a vacuum counterterm rather than a pure current.  Any nonlinear
\(F\) is a different current EFT and must pass the same closure and perturbative
null tests rather than being identified with the minimal branch.
\end{criterion}

Equivalently, the stable barotropic current class may be generated by a function
\begin{equation}
0\le s(n)\le1,\qquad
s(n):=\frac{\dd\ln F'(n)}{\dd\ln n},
\end{equation}
which gives
\begin{align}
F'(n)&=F'(n_0)
\exp\!\left[\int_{n_0}^{n}s(x)\,\dd\ln x\right],\\
F(n)&=F(n_0)+\int_{n_0}^{n}F'(y)\,\dd y .
\end{align}
The case $s(n)=0$ gives $F(n)=\mu c^2n+\mathrm{const}$, i.e. dust plus a
vacuum offset; the offset must be assigned to the vacuum sector to avoid double
counting with the four-form branch.  The case $s(n)=w$ with $0<w\le1$ gives
$F(n)\propto n^{1+w}$ and $p=w\eps$; it is stable but not the minimal dust
residual.

To use the degree as the Brown-current number one must impose the sector
identification
\begin{equation}
N_J:=\int_{\Sigma_a}n\,\dd V_a=|Q|.
\label{eq:degree_current_identification_v14}
\end{equation}
This is an additional branch assignment, not a consequence of the degree
definition alone: \(Q\in\mathbb Z\) labels a topological sector, and only the
extra branch rule \(N_J=|Q|\) turns it into a Brown-current number.  With this
assignment and homogeneous coarse graining on
\(\Sigma_{A_0}=S^3_{A_0}/\Gamma\),
\begin{equation}
\Vol(\Sigma_{A_0})=\frac{2\pi^2A_0^3}{|\Gamma|},
\qquad
n(a)=\frac{|Q|}{a^3\Vol(\Sigma_{A_0})},
\label{eq:volume_number_density_v5}
\end{equation}
and therefore
\begin{equation}
\Ow=
\frac{\mu_{\wind}c^2|Q|}{\Vol(\Sigma_{A_0})\eps_{\crit,0}}
=
\frac{4G\mu_{\wind}|Q|\,|\Gamma|}{3\pi H_0^2A_0^3} .
\label{eq:omega_w_v5}
\end{equation}
This is an exact bookkeeping identity only after the branch postulate
\(N_J=|Q|\) and the non-topological data \((\mu_{\wind},|\Gamma|,A_0)\) are
supplied.  It is not an amplitude prediction.

Four-form vacuum sectors of this kind are standard in gravitational theory~\cite{Aurilia,HenneauxTeitelboim,BrownTeitelboim,BoussoPolchinski}.  For a four-form \(F_4=\dd C_3\), take
\begin{equation}
S_F=-\frac{Z_F}{2\cdot4!}\int \sqrt{-g}\,
F_{\mu\nu\rho\sigma}F^{\mu\nu\rho\sigma}\,\dd^4x,
\qquad Z_F>0 .
\label{eq:four_form_action_v5}
\end{equation}
The stress tensor obtained from Eq.~\eqref{eq:four_form_action_v5} is
\begin{equation}
T^F_{\mu\nu}
=
\frac{Z_F}{3!}F_{\mu\alpha\beta\gamma}
F_{\nu}{}^{\alpha\beta\gamma}
-
\frac{Z_F}{2\cdot4!}g_{\mu\nu}F_{\alpha\beta\gamma\delta}
F^{\alpha\beta\gamma\delta} .
\label{eq:four_form_stress_general_v7}
\end{equation}
The field equation implies that the four-form is proportional to the metric
volume form on each connected branch,
\begin{equation}
F_{\mu\nu\rho\sigma}=f\varepsilon_{\mu\nu\rho\sigma},
\qquad
f=\text{constant},
\label{eq:four_form_branch_v8}
\end{equation}
up to the orientation convention for the Hodge dual.  Only \(f^2\) enters the
stress tensor.  Define
\begin{equation}
\qhat=\sqrt{Z_F}f,
\qquad
\eps_q=\frac12\qhat^2 .
\label{eq:four_form_energy_v5}
\end{equation}
With the \((-+++)\) signature,
\begin{equation}
\varepsilon_{\mu\alpha\beta\gamma}
\varepsilon_{\nu}{}^{\alpha\beta\gamma}=-3!g_{\mu\nu},
\qquad
F_{\alpha\beta\gamma\delta}F^{\alpha\beta\gamma\delta}=-4!f^2,
\end{equation}
which gives
\begin{equation}
T^q_{\mu\nu}=-\eps_q g_{\mu\nu},
\qquad
\Lambda_{\eff}=\frac{8\pi G}{c^4}\eps_q .
\label{eq:four_form_stress_v5}
\end{equation}

\begin{proposition}[Three separate four-form statements]
\textup{For the four-form model considered here, three logically distinct statements must not be conflated.}
First,
\begin{equation}
F_4=f\,\varepsilon_4
\quad\Rightarrow\quad
T^q_{\mu\nu}=-\eps_q g_{\mu\nu},
\qquad
\eps_q=\frac12\qhat^2 .
\label{eq:four_form_local_branch_v23}
\end{equation}
This is local stress-tensor closure on a connected fixed branch and follows from
Eqs.~\eqref{eq:four_form_stress_general_v7}--\eqref{eq:four_form_stress_v5}.
Second,
\begin{equation}
\eps_q=\frac12\qhat^2,
\qquad
\qhat\in\R,
\label{eq:four_form_continuous_branch_v23}
\end{equation}
is a continuous amplitude statement.  Third,
\begin{equation}
\qhat=\NF\eF,
\qquad
\NF\in\Z,
\label{eq:four_form_quantization_rule_v23}
\end{equation}
is an additional flux-quantization rule.  The third statement is not implied by
compact spatial topology alone; it requires a charge normalization, compact
zero-mode normalization, membrane sector, or equivalent branch rule.
\end{proposition}
For a fixed superselection branch, the mixed-index perturbation vanishes,
\begin{equation}
\delta T^{q\,\mu}{}_{\nu}=0,
\label{eq:mixed_index_four_form_v5}
\end{equation}
while the covariant-index variation
\begin{equation}
\delta T^q_{\mu\nu}=-\eps_q\delta g_{\mu\nu}
\label{eq:covariant_four_form_v5}
\end{equation}
is purely metric-induced.

\begin{proposition}[Closure versus flux-lattice selection]
The homogeneous closure criterion determines only the stress-energy form
\begin{equation}
T^q_{\mu\nu}=-\eps_q g_{\mu\nu},
\qquad
\eps_q=\text{constant}\ge0 .
\end{equation}
It does not by itself determine whether the constant belongs to a discrete flux
lattice.  Spatial compactness by itself does not quantize the four-form
amplitude; the lattice enters only after a charge normalization, compact
zero-mode normalization, or membrane sector has been specified.  If such a
normalization is supplied so that
\begin{equation}
\qhat=\NF\eF,
\qquad
\NF\in\Z,
\end{equation}
then the admissible vacuum amplitudes are restricted to
\begin{equation}
\eps_q=\frac12\NF^2\eF^2 .
\label{eq:flux_lattice_v5}
\end{equation}
Thus the FLRW closure theorem is a continuous stress-energy classification,
whereas flux quantization is an additional amplitude-selection constraint.
Arbitrary residual matching requires an unfixed normalization \(\eF\), a dense
landscape of allowed branches, or a phenomenological effective constant.
\end{proposition}

For an observational target $\eps_q^{\rm target}$ and a single flux unit $\eF$,
the corresponding branch selector is the lattice problem
\begin{equation}
\NF^\star\in
\operatorname*{argmin}_{\NF\in\Z}
\left|\frac{1}{2}\NF^2\eF^2-\eps_q^{\rm target}\right| .
\end{equation}
Exact fixed-vacuum closure requires
\begin{equation}
\frac{1}{2}(\NF^\star)^2\eF^2=\eps_q^{\rm target},
\end{equation}
whereas phenomenological closure only requires the mismatch to lie below the
adopted observational tolerance.  For the fixed-four-form model considered
here, the amplitude routes are a continuous EFT branch $q_b\in\R$, a single
discrete flux $q_b=\NF\eF$, a dense multiflux lattice
$\eps_q=\frac{1}{2}\sum_I N_I^2e_I^2$, a directly fitted phenomenological
constant, or no admissible fixed branch.


\begin{criterion}[Fixed four-form branch]
On a scale-factor interval \(I\), the fixed-vacuum source law requires
\begin{equation}
\begin{gathered}
\frac{\dd \NF}{\dd\ln a}=0,
\qquad
\frac{\dd \eF}{\dd\ln a}=0,\\
\text{no membrane transition on }I .
\end{gathered}
\label{eq:fixed_branch_v13}
\end{equation}
This can be imposed either as strict superselection,
\begin{equation}
\mathcal H=\bigoplus_{\NF}\mathcal H_{\NF},
\qquad
|\Psi_{\rm cos}\rangle\in\mathcal H_{\NF^\star},
\end{equation}
or as EFT transition suppression,
\begin{equation}
\Gamma_{\rm mem}\,\mathcal V_4(I)\ll1 .
\end{equation}
If membrane transitions are allowed without suppression in a flux-lattice
branch, then
\begin{equation}
\NF\longrightarrow \NF\pm n_{\rm mem},
\qquad n_{\rm mem}\in\Z,
\end{equation}
so that
\begin{equation}
\eps_q:\frac{1}{2}\NF^2\eF^2
\longrightarrow
\frac{1}{2}(\NF\pm n_{\rm mem})^2\eF^2,
\end{equation}
and the theorem applies only piecewise between jumps; the residual is not
globally fixed vacuum.
\end{criterion}


\section{Minimal background system and compact causal ledger}

The minimal background system represented by the closed homogeneous atlas is
\begin{equation}
G_{\mu\nu}=\frac{8\pi G}{c^4}
\left(T^{\vis}_{\mu\nu}+T^{\wind}_{\mu\nu}+T^q_{\mu\nu}\right),
\label{eq:minimal_background_system_v25}
\end{equation}
with
\begin{equation}
T^{\wind}_{\mu\nu}=\eps_{\wind}\frac{U_\mu U_\nu}{c^2},
\qquad
\eps_{\wind}=\mu_{\wind}c^2 n,
\qquad
\nabla_\mu(nU^\mu/c)=0,
\label{eq:minimal_dust_system_v25}
\end{equation}
and
\begin{equation}
T^q_{\mu\nu}=-\eps_q g_{\mu\nu}.
\label{eq:minimal_vacuum_system_v25}
\end{equation}
For the closed round cover this gives the background Friedmann relation
\begin{equation}
H^2+\frac{c^2}{a^2A_0^2}
=
\frac{8\pi G}{3c^2}
\left(\eps_{\vis}+\eps_{\wind,0}a^{-3}+\eps_q\right).
\label{eq:minimal_friedmann_v25}
\end{equation}
Equation~\eqref{eq:minimal_friedmann_v25} is not an amplitude derivation.  It is
the background equation after the constants \((\eps_{\wind,0},\eps_q,A_0)\) and
the visible sector have been supplied.

The compact causal part of the framework is source-free.  On the simply
connected closed cover
\begin{equation}
\dd s^2=-c^2\dd t^2+a^2(t)A_0^2\dd\Omega_3^2,
\label{eq:causal_closed_metric_v25}
\end{equation}
null propagation defines the dimensionless conformal causal distance
\begin{equation}
\Delta\chi(t_1,t_2)=\int_{t_1}^{t_2}\frac{c\,\dd t}{a(t)A_0}.
\label{eq:delta_chi_v25}
\end{equation}
Because the geodesic diameter of the unit three-sphere is \(\pi\), full causal
reach on the cover requires
\begin{equation}
\Delta\chi(t_1,t_2)\ge\pi .
\label{eq:full_causal_reach_v25}
\end{equation}
For a static cover \(a=1\), this reduces to
\begin{equation}
t_{\rm coh}=\frac{\pi A_0}{c}.
\label{eq:static_coherence_time_v25}
\end{equation}
The corresponding cover-volume accessibility ratio is
\begin{equation}
\begin{gathered}
\mathcal G_{\rm cau}(t_1,t_2)=\frac{\Vol B_{\chi}}{\Vol S^3},\\
\mathcal G_{\rm cau}(t_1,t_2)=
\frac{\chi-\frac12\sin(2\chi)}{\pi},
\qquad
\chi=\min\{\Delta\chi,\pi\}.
\end{gathered}
\label{eq:causal_access_ratio_v25}
\end{equation}
This is the invariant content of the compact causal ledger on the simply
connected cover: light supplies null cones, the compact cover supplies a finite
causal diameter, and expansion enters only through conformal time.  The causal
ledger does not modify Eq.~\eqref{eq:minimal_background_system_v25} and does not
select \(\eps_{\wind,0}\) or \(\eps_q\).

For a quotient \(\Sigma_\Gamma=S^3/\Gamma\), Eq.~\eqref{eq:delta_chi_v25}
remains the local null-reach variable on the cover, but the exact accessibility
condition is quotient-relative.  For an observer \(x_o\), define
\begin{equation}
D_\Gamma(x_o):=\sup_{y\in\Sigma_\Gamma}d_\Gamma(x_o,y),
\label{eq:quotient_causal_radius_v26}
\end{equation}
where \(d_\Gamma\) is the quotient geodesic distance.  Full causal access to the
quotient from \(x_o\) is equivalent to
\begin{equation}
\Delta\chi(t_1,t_2)\ge D_\Gamma(x_o),
\label{eq:quotient_full_access_v26}
\end{equation}
and the invariant accessibility fraction is
\begin{equation}
\mathcal G^\Gamma_{\rm cau}(x_o;t_1,t_2)
=
\frac{\Vol\{y\in\Sigma_\Gamma:d_\Gamma(x_o,y)\le\Delta\chi\}}
{\Vol(\Sigma_\Gamma)} .
\label{eq:quotient_causal_fraction_v26}
\end{equation}
For the simply connected cover, \(D_\Gamma=\pi\) and
Eq.~\eqref{eq:quotient_causal_fraction_v26} reduces to
Eq.~\eqref{eq:causal_access_ratio_v25}.  Thus compact causal closure requires
the same data as compact spectral closure: the quotient \(\Gamma\), its
projector, and the observer-position prior.  A model with no specified topology
is locally viable but globally incomplete.

\begin{criterion}[Causal ledger is not a source]
Compact causal accessibility is a propagation and global-identification
criterion.  It may constrain which regions or repeated images are mutually
causally accessible, and it may enter topology-covariance calculations through
the quotient geometry.  It is not a local homogeneous stress tensor and cannot
replace the Brown-current mass scale, the four-form flux unit, or the compact
Boltzmann transfer map.
\end{criterion}

\begin{proposition}[Local Euclidean limit]
Let \(L\) be the physical size of a laboratory or survey patch with
\(L\ll aA_0\).  In normal coordinates on the spatial slice, the compact
spherical metric differs from the Euclidean FLRW metric by relative corrections
of order \(L^2/(a^2A_0^2)\).  Therefore the residual atlas, Brown-current dust
law, and fixed four-form vacuum law are locally the same EFT statements as in an
\(R^3\) FLRW chart.  Compactness first appears through global identifications,
spectra, causal reach, and covariance, not through a different local stress
tensor.
\end{proposition}

\section{Compact topology and quotient projection}

The spatial topology is a global input, not a local equation of state.
The local residual classification above did not require compactness or positive
curvature.  The global ledger in this section does: throughout the
topological-current and compact three-form slots we restrict to compact
orientable spherical space forms, equivalently to finite freely acting
orientation-preserving isometry groups \(\Gamma\) on \(S^3\).  Without
orientability the degree and top-cohomology statements must be replaced by their
twisted-coefficient versions; such branches are not considered here.  The
classification of complete constant-curvature space forms gives~\cite{HawkingEllis,Wolf}
\begin{equation}
\Sigma=\widetilde\Sigma/\Gamma,
\qquad
\widetilde\Sigma\in\{S^3,\R^3,H^3\} .
\end{equation}
In the compact positive-curvature case,
\begin{equation}
\Sigma=S_A^3/\Gamma,
\qquad
\Vol(\Sigma)=\frac{2\pi^2A^3}{|\Gamma|} .
\label{eq:space_form_volume_v5}
\end{equation}
If simple connectivity is assumed, \(\Gamma=\id\) and \(\Sigma\simeq S_A^3\).
Without simple connectivity, the invariant conclusion is only that the
universal cover is \(S_A^3\).

On the simply connected cover,
\begin{equation}
-\Delta_{S_A^3}Y_{n\ell m}=\frac{n(n+2)}{A^2}Y_{n\ell m},
\qquad
{\rm mult}(n)=(n+1)^2 .
\label{eq:s3_spectrum_v5}
\end{equation}
For a quotient, fields are \(\Gamma\)-invariant cover fields.  The projector is
\begin{equation}
\boxed{
\PG=\frac1{|\Gamma|}\sum_{\gamma\in\Gamma}\gamma^* .
}
\label{eq:quotient_projector_v5}
\end{equation}
The physical scalar eigenspace is \(\Hn^\Gamma=\PG\Hn\), with multiplicity
\begin{equation}
m_\Gamma(n)=\Tr_{\Hn}\PG .
\label{eq:quotient_multiplicity_v5}
\end{equation}
The same logic applies to physical vector and tensor harmonic sectors after
closed-FLRW gauge and constraint modes are removed.

Compact topology can therefore enter through
\begin{equation}
\Vol(\Sigma),\qquad \Spec_\Gamma(\Delta),\qquad m_\Gamma(n),\qquad
\mathrm{Cov}_\Gamma,
\label{eq:topology_channels_v5}
\end{equation}
not through a new local homogeneous stress tensor.  Observer-position dependence
appears for non-globally-homogeneous quotients; for globally homogeneous
quotients the covariance is still quotient-projected but is not position
dependent in the same way.

\begin{criterion}[Topology-evidence boundary]
A compact spherical model may claim compact-background consistency after
specifying \(S^3/\Gamma\) and the quotient projection.  It may not claim topology
evidence or benchmark superiority until the projector has been propagated into
transfer functions and observable covariance,
\begin{equation}
\PG\longrightarrow
m_\Gamma(n),\ \Xi_n^\Gamma,\
\Delta_{\ell n}^{X,\Gamma},\ C_\ell^{XY,\Gamma},\
\mathrm{Cov}_\Gamma,\ \mathcal L_\Gamma .
\label{eq:transfer_boundary_v5}
\end{equation}
A background quotient without this propagated covariance is a compact background
ansatz, not a topology likelihood.
\end{criterion}

\section{Amplitude-selection boundary}

Compact topology supplies labels, volume, and spectral restrictions.  It does
not supply numerical cosmological abundances.  In particular,
\begin{equation}
(\Gamma,[\Sigma,S^3],H^3(\Sigma,\Z),Q,\NF)
\not\Rightarrow
(\Ow,\Oq,\Lambda_{\eff}) .
\label{eq:not_amplitudes_v5}
\end{equation}
Even after the geometric radius \(A_0\) is supplied, the amplitudes require
\(\mu_{\wind}\), \(\eF\), branch labels, compact-gauge normalization, and a
selection rule or phenomenological prior.  This is an amplitude-selection
boundary, not a failure of the closure theorem.


\begin{theorem}[Topological decoupling of amplitudes]
\label{thm:topological_decoupling_v12}
Fix the class of compact orientable spherical space forms and the discrete
sector data
\[
\mathcal T_{\rm top}=(\Gamma,[\Sigma,S^3],H^3(\Sigma,\Z),Q,\NF) .
\]
There is no convention-independent, dimensionally complete topology-only map
from \(\mathcal T_{\rm top}\) to the physical amplitude pair
\((\Ow,\Oq)\).  Any physical amplitude selector must factor through additional
non-topological data,
\begin{equation}
\mathfrak S_{\rm amp}:
(\mathcal T_{\rm top};A_0,H_0,\mu_{\wind},\eF,\mathcal B)
\longrightarrow
(\Ow,\Oq),
\label{eq:split_amp_selector_v12}
\end{equation}
where \(\mathcal B\) denotes branch, compact-gauge, boundary, or prior data.
Equivalently, the microscopic selector splits into independent current and
vacuum slots,
\begin{equation}
\mathcal K_{\rm micro}=\mathcal K_{\wind}\oplus\mathcal K_q,
\qquad
\mathcal K_{\wind}\ni \mu_{\wind},
\qquad
\mathcal K_q\ni \eF .
\label{eq:micro_selector_split_v12}
\end{equation}
Thus topology fixes admissible sectors and projected spectra, whereas
microphysics or phenomenological branch selection fixes amplitudes.
\end{theorem}

\begin{proof}
Holding \((\Gamma,Q,\NF)\) fixed while varying \(\mu_{\wind}\), \(\eF\),
\(A_0\), or \(H_0\) changes \((\Ow,\Oq)\).  Hence the amplitudes cannot be
functions of topology alone.  Equivalently, the entries in
\(\mathcal T_{\rm top}\) label sectors and restrict mode spaces, but they contain
no conversion from charge number to energy density.  For the current branch,
\begin{equation}
\Ow=
\frac{\mu_{\wind}c^2|Q|}{\Vol(\Sigma_{A_0})\eps_{\crit,0}},
\qquad
\Vol(\Sigma_{A_0})=\frac{2\pi^2A_0^3}{|\Gamma|} .
\label{eq:current_amp_decoupling_v12}
\end{equation}
Therefore \(Q\), \(\Gamma\), and the topological class do not determine
\(\Ow\) unless the effective mass scale \(\mu_{\wind}\), the geometric scale
\(A_0\), and the critical-density normalization are supplied.  For the
four-form branch,
\begin{equation}
\Oq=\frac{\eps_q}{\eps_{\crit,0}},
\qquad
\eps_q=\frac12\NF^2\eF^2 .
\label{eq:four_form_amp_decoupling_v12}
\end{equation}
Therefore the flux integer \(\NF\) does not determine \(\Oq\) unless the flux
unit \(\eF\) and the same cosmological normalization are supplied.
Consequently, any numerical map from topological labels alone to \((\Ow,\Oq)\)
is either a phenomenological assignment or a hidden use of additional
microphysical or boundary data.  This proves the decoupling statement.
\end{proof}

A useful diagnostic for selector strength is the curvature-linear ansatz
\(\Ow=C|\Omega_K|\).  If observations impose \(|\Omega_K|<\kappa\) and the
target is \(\Omega_\star\), then
\begin{equation}
C>\frac{\Omega_\star}{\kappa} .
\label{eq:selector_bound_v5}
\end{equation}
For \(\Omega_\star\simeq0.26\) and \(\kappa=0.005\), one needs \(C>52\).
A toy holographic winding saturation condition,
\begin{equation}
\mu_{\wind}c^2|Q|=\frac{c^4A_0}{2G},
\end{equation}
gives only
\begin{equation}
\Ow^{\rm HWS}=\frac{2}{3\pi}|\Omega_K|
\label{eq:hws_simple_v5}
\end{equation}
for the simply connected case.  For a quotient, imposing the same total-energy
saturation on the quotient volume gives
\begin{equation}
\Ow^{\rm HWS}(\Gamma)=\frac{2|\Gamma|}{3\pi}|\Omega_K| .
\label{eq:hws_quotient_v5}
\end{equation}
This is not a derivation of \(\Omega_c\).  If \(|\Gamma|\sim 10^2\) is used to
obtain a large effective coefficient, the proposal becomes a large-quotient
amplitude ansatz that must pass matched-circle, observer-position, covariance,
and spectral-multiplicity checks.

\section{Perturbative and observational null tests}

We use standard gauge-invariant cosmological perturbation conventions~\cite{Bardeen,KodamaSasaki,Mukhanov,MaBertschinger}.  At linear order, the residual is defined before gauge choice by
\begin{equation}
\delta T^\D_{\mu\nu}=\frac{c^4}{8\pi G}\delta G_{\mu\nu}
-\delta T^{\vis}_{\mu\nu} .
\label{eq:linear_residual_v5}
\end{equation}
The fixed four-form branch contributes no independent scalar perturbation on a
fixed branch.  After subtracting this fixed vacuum component, the minimal
current branch must obey the dust null conditions on every physical
quotient-projected scalar mode:
\begin{equation}
\boxed{
\delta p_{\wind,{\rm rf}}=0,
\qquad
\Pi_{\wind}^{\GI}=0,
\qquad
\frac{v_{s,\wind}^2}{c^2}=0 .
}
\label{eq:dust_null_conditions_v5}
\end{equation}
A statistically significant intrinsic pressure perturbation, rest-frame sound
speed, anisotropic stress, entropy source, or independent non-dust stress source
not assigned to the visible sector falsifies the minimal perturbative branch.

Current background data sharpen the closure tests without changing their
logical status.  Planck 2018 combined with BAO gives a curvature constraint
consistent with flatness, conventionally quoted as
\(\Omega_K=0.0007\pm0.0019\) for TT,TE,EE+lowE+lensing+BAO~\cite{Planck2018}.
Recent DESI DR2 BAO results are well described by flat \(\Lambda\)CDM, while
BAO+CMB prefers a \(w_0w_a\) model over \(\Lambda\)CDM at \(3.1\sigma\), and
BAO+CMB plus supernova samples gives a sample-dependent preference in the range
\(2.8\)--\(4.2\sigma\)~\cite{DESIDR2}.  If a reconstructed residual
pressure evolution remains significant after systematics, dataset tension, and
visible-sector convention dependence are accounted for, then the fixed-four-form
minimal branch fails the
\(\Rcl_q=0\) null test.  That would falsify the fixed-vacuum source law,
not compact spherical topology.

Likewise, matched-circle and compact-covariance limits constrain detectable
domains and topology families, not every compact spherical quotient.  The
published lens-space constraints are strong statements for specified
\(L(p,q)\), curvature-scale ranges, and circle-search domains; they are not a
completed scan of all \(S^3/\Gamma\) likelihoods.  Lens-space constraints are
therefore constraints on specified \(S^3/\Gamma\) implementations and their
observer-position/covariance pipelines, where applicable~\cite{Cornish,
PlanckTopology,SahaLens}.

This separation also appears in compact-topology covariance studies.  The
quotient changes the allowed eigenmodes and covariance.  Assumptions about
large-scale primordial power can amplify or suppress detectability of the
quotient signal~\cite{COMPACTIIa,COMPACTIIb,COMPACTIIIb,COMPACTIIc}.  A
Casimir stress is a different physical source, not an amplitude selector for the
minimal homogeneous branch.  If it produces anisotropic backreaction, it exits
the homogeneous locally-FLRW residual class rather than selecting the amplitudes
inside it~\cite{NegroCasimir}.


\section{What would be required for a topology-likelihood claim}
\label{sec:transfer_protocol_v16}

The following section is not part of the homogeneous closure theorem.  It records
the additional ingredients required before a compact-background ansatz could be
promoted to a topology-evidence calculation.  The homogeneous theorem is a
background classification: a topology-evidence claim begins only after the
quotient projector is propagated to the covariance of an actual data vector.
The admissible completion is therefore a protocol, not a detection claim.  It
requires a concrete quotient $\Gamma$, a data vector, a mask, a foreground and
noise model, a compact Boltzmann implementation, and a likelihood comparison.

\subsection{SVT harmonic ledger}

Let $r=s,v,t$ denote scalar, transverse-vector, and transverse-traceless tensor
sectors after gauge and constraint modes are removed.  Quotient harmonics are
obtained by projection from the cover,
\begin{equation}
\PG=\frac{1}{|\Gamma|}\sum_{\gamma\in\Gamma}\gamma^*,
\qquad
Y^\Gamma_{n\alpha}=\PG Y_{n\alpha},
\end{equation}
and analogously for $V^\Gamma_{n\alpha i}$ and $T^\Gamma_{n\alpha ij}$.
Because scalar, vector, and tensor spectra need not share the same eigenvalue
formula, write
\begin{equation}
-\Delta_r Q^{(r),\Gamma}_{n\alpha}
=\lambda_n^{(r)} Q^{(r),\Gamma}_{n\alpha},
\qquad r=s,v,t .
\end{equation}
The quotient multiplicities are
\begin{equation}
m_\Gamma^{(r)}(n)
=
\dim\ker\bigl(\Delta_r+\lambda_n^{(r)}\bigr)^\Gamma,
\qquad r=s,v,t .
\label{eq:sector_multiplicity_v16}
\end{equation}
Here the superscript $\Gamma$ means the invariant subspace selected by $\PG$.

In Newtonian gauge,
\begin{equation}
\dd s^2=a^2(\eta)\{-(1+2\Psi)\dd\eta^2
+(1-2\Phi)\gamma_{ij}\dd x^i\dd x^j\} .
\end{equation}
For the minimal dust current one may write, up to the sign convention for
scalar harmonics and metric potentials,
\begin{align}
\delta'_{\wind,n\alpha}&=-\theta_{\wind,n\alpha}
+3\Phi'_{n\alpha},\\
\theta'_{\wind,n\alpha}+\mathcal H\theta_{\wind,n\alpha}
&=\lambda_n^{(s)}\Psi_{n\alpha} .
\end{align}
The necessary null requirements are
\begin{equation}
\delta p_{\wind,{\rm rf}}=0,
\qquad
\Pi_\wind^{\GI}=0,
\qquad
v_{s,\wind}^2=0 .
\label{eq:pert_null_tests_v16}
\end{equation}
For a stable non-dust current EFT, a schematic scalar ledger is
\begin{align}
\delta'_{J,n\alpha}
&=-(1+w_J)(\theta_{J,n\alpha}-3\Phi'_{n\alpha})
-3\mathcal H(c_s^2-w_J)\delta_{J,n\alpha},\\
\theta'_{J,n\alpha}
&=-\mathcal H(1-3c_s^2)\theta_{J,n\alpha}
-\frac{w'_J}{1+w_J}\theta_{J,n\alpha}\\
&\quad+\frac{c_s^2}{1+w_J}\lambda_n^{(s)}\delta_{J,n\alpha}
+\lambda_n^{(s)}\Psi_{n\alpha} .
\end{align}
Vector and tensor null conditions for the residual are
\begin{align}
\Pi^{(v)}_{\wind,n\alpha}&=0,&
 v^{(v)\prime}_{\wind,n\alpha}+\mathcal H v^{(v)}_{\wind,n\alpha}&=0,\\
\Pi^{(t)}_{\wind,n\alpha}&=0,&
\Pi^{(t)}_{q,n\alpha}&=0 .
\end{align}
For tensor perturbations, use the kinetic operator actually appearing in the
linearized Einstein equation on the compact quotient,
\begin{equation}
\mathcal K_t Q^{(t),\Gamma}_{n\alpha}
=\kappa^{(t)}_n Q^{(t),\Gamma}_{n\alpha} .
\end{equation}
This convention-safe notation absorbs the curvature shift that would otherwise
depend on whether one calls the eigenvalue a rough-Laplacian, Lichnerowicz, or
already-shifted kinetic eigenvalue.  The tensor amplitudes obey schematically
\begin{equation}
h''_{n\alpha}+2\mathcal Hh'_{n\alpha}
+\kappa^{(t)}_n h_{n\alpha}
=\frac{16\pi G}{c^4}a^2\Pi^{(t)}_{{\rm vis},n\alpha} .
\end{equation}
For the four-form sector, on a fixed branch,
\begin{equation}
\delta T^{q\,\mu}{}_{\nu}=0,
\qquad
\delta T^q_{\mu\nu}=-\eps_q\delta g_{\mu\nu} .
\end{equation}
Thus the branch is unperturbed as a physical vacuum component; the covariant
index variation is metric-induced.

\subsection{Compact transfer functions and covariance}

For an observable $X\in\{T,E,B,\phi_{\rm lens},g,\ldots\}$ define
\begin{equation}
a^X_{\ell m}(x_o,R)=
\sum_{n,\alpha,I}
\Delta^{X,\Gamma,I}_{\ell m;n\alpha}(x_o,R;\eta)
\mathcal R^I_{n\alpha} .
\end{equation}
Here $x_o$ is the observer position in $\Sigma_\Gamma$, $R\in SO(3)$ is sky
orientation, and $I$ labels scalar, vector, tensor, or other declared initial
sectors.  The quotient-projected primordial covariance is
\begin{equation}
\langle \mathcal R^I_{n\alpha}\mathcal R^{J*}_{n'\beta}\rangle
=
\Xi^{IJ,\Gamma}_{n\alpha,n'\beta} .
\end{equation}
Hence, using multi-indices $A=(n,\alpha,I)$ and $B=(n',\beta,J)$,
\begin{equation}
{\rm Cov}^{XY,\Gamma}_{\ell m,\ell' m'}(x_o,R)
=
\sum_{A,B}
\Delta^{X,\Gamma}_{\ell m;A}\,
\Xi^{\Gamma}_{AB}\,
\Delta^{Y,\Gamma*}_{\ell' m';B} .
\label{eq:compact_covariance_v16}
\end{equation}
The required chain is
\begin{equation}
\PG\longrightarrow
m_\Gamma^{(r)}(n)\longrightarrow
\Xi^\Gamma\longrightarrow
\Delta^{X,\Gamma}\longrightarrow
{\rm Cov}^\Gamma\longrightarrow
\mathcal L_\Gamma .
\label{eq:transfer_chain_v16}
\end{equation}
At the perturbative level, the exact degeneracy statement is
\begin{equation}
\begin{aligned}
&\text{same }(\eps_{\wind,0},\eps_q),
\quad \text{same visible sector},\\
&\text{same }\Xi^\Gamma,
\quad \text{same }\Delta^{X,\Gamma}
\Longrightarrow
\text{same }{\rm Cov}^{XY,\Gamma},\\
&\Longrightarrow
\text{no observable distinction from CDM}+\Lambda .
\end{aligned}
\label{eq:perturbative_degeneracy_v23}
\end{equation}
within the declared observable data model.  Initial statistics and transfer
functions determine the covariance; the covariance determines the likelihood.  A
background quotient without the propagated covariance in Eq.~\eqref{eq:transfer_chain_v16}
is a compact background ansatz, not a topology likelihood.

\subsection{Observer-position priors and likelihood}

For globally homogeneous quotients the observer-position integral collapses to
one orbit.  For non-globally-homogeneous quotients one must integrate over the
inequivalent observer domain
\begin{equation}
\mathcal F_\Gamma=\Sigma_\Gamma/{\rm Isom}(\Sigma_\Gamma) .
\end{equation}
The default invariant prior is volume-uniform,
\begin{equation}
\pi_\Gamma(x_o)=\frac{1}{\Vol(\mathcal F_\Gamma)},
\qquad
\int_{\mathcal F_\Gamma}\pi_\Gamma\,\dd\mu_\Gamma=1 .
\end{equation}
More general priors may be used only if declared before fitting,
\begin{equation}
\pi_\Gamma(x_o|\mathcal S)\propto
\pi_{\rm geom}(x_o)W_{\rm selection}(x_o) .
\end{equation}
Hidden observer tuning is not an admissible topology claim.

For a fixed $(\Gamma,x_o,R,\eta)$ the masked data covariance is written as
\begin{equation}
C_\Gamma=M\,{\rm Cov}_\Gamma(x_o,R,\eta)\,M^\dagger+N+C_{\rm fg}(\nu) .
\end{equation}
Here $M$ is the mask or mode-coupling operator, $N$ is instrumental noise, and
$\nu$ denotes declared foreground nuisance parameters.  The covariance-level
Gaussian likelihood is
\begin{align}
\mathcal L_\Gamma(d)
&=
\int_{\mathcal F_\Gamma}\pi_\Gamma(x_o)\,\dd\mu_\Gamma(x_o)
\int_{SO(3)}\dd R
\nonumber\\
&\quad\times
\int\dd\eta\,\dd\nu\;
\mathcal G(d;C_\Gamma)\Pi(\eta,\nu),
\label{eq:likelihood_v16}\\
\mathcal G(d;C)&=
\frac{\exp[-\frac{1}{2}d^\dagger C^{-1}d]}{\sqrt{\det(2\pi C)}} .
\end{align}
The comparison model must use the same local expansion history, curvature,
visible sector, background amplitudes, mask, noise, and foreground conventions,
differing only in the simply connected or statistically isotropic covariance
benchmark.  We denote this matched benchmark by \(\mathcal L_{\rm bench}\).  The
first point at which topology can become a likelihood claim is
\begin{equation}
\Delta\ln\mathcal L=
\ln\mathcal L_\Gamma-
\ln\mathcal L_{\rm bench} .
\end{equation}

\subsection{Transfer diagnostics}
Blind-injection recovery, matched-circle rejection, and low-\(\ell\) covariance comparisons are necessary implementation diagnostics for any chosen quotient.  They are not part of the homogeneous classification theorem and are stated explicitly in Appendix~\ref{app:transfer_diagnostics}.

\subsection{Closure-route ledger}

After the preceding definitions, the remaining solution space is finite in
structure.  In the current sector, pure pressureless behavior requires
$F(n)=\mu c^2n$ after any additive vacuum counterterm has been removed and
assigned consistently to the vacuum sector; nonlinear stable $F$ defines a
different EFT unless its pressure and sound speed are negligible.  In the
fixed-four-form sector considered here,
the amplitude routes are continuous $q_b$, a single flux lattice, a dense
multiflux lattice, a phenomenological constant, or no branch.  In the topology
sector, topology can enter only through
\begin{equation}
\Vol(\Sigma_\Gamma),\quad
\PG,
\quad m_\Gamma^{(r)}(n),
\quad \Xi^\Gamma,
\quad {\rm Cov}^\Gamma,
\quad \pi_\Gamma(x_o) .
\end{equation}
It cannot determine $\Omega_c$, $\Omega_\Lambda$, or $\Lambda_{\eff}$ without
$\mu_{\wind}$, $\eF$, $A_0$, $H_0$, and branch data.

\begin{theorem}[Closure-route ledger]
Within the declared FLRW residual, Brown-current, fixed-four-form, and
quotient-covariance ledgers, every strict extension of the minimal homogeneous
residual enters through at least one of the following routes:
\begin{enumerate}
\item \(\Rcl_q\ne0\) or \(\Rcl_{\wind}\ne0\):
a non-minimal homogeneous equation of state.
\item \(\delta p_{\wind,{\rm rf}}\ne0\), \(\Pi_{\wind}^{\GI}\ne0\), or
\(v_{s,\wind}^2\ne0\): non-dust perturbations.
\item \(\dd \NF/\dd\ln a\ne0\), \(\dd \eF/\dd\ln a\ne0\), or membrane
transitions: a non-fixed vacuum branch.
\item \(\mathfrak S_{\rm amp}\) uses non-topological normalization or branch
data: amplitude selection.
\item \({\rm Cov}^\Gamma\), \(\pi_\Gamma(x_o)\), \(M\), \(N\), and
\(C_{\rm fg}\) enter a likelihood: topology evidence.
\end{enumerate}
Thus a manuscript-level claim is invariantly classified as a background closure
claim, a perturbative source claim, a microscopic branch-selection claim, or a
topology-likelihood claim.
\end{theorem}

\begin{proof}
The atlas \((\Aq,\Aw)\) is algebraically invertible, so any homogeneous
departure from the minimal residual is exactly a departure from the constant
nonnegative submanifold, giving (i).  The Brown-current representative is dust
only on the linear branch \(F(n)=\mu c^2n\) after vacuum offsets are assigned to
the vacuum sector, giving (ii).  A four-form is a fixed vacuum source only on a
fixed branch with no membrane transition, giving (iii).  The topological
decoupling theorem proves that amplitudes require non-topological
normalization, giving (iv).  Finally, the quotient projector becomes an
observational topology claim only after it has been propagated to a masked,
noisy data covariance and likelihood, giving (v).  These cases exhaust only the
objects introduced in the residual, source, topology, and likelihood ledgers.
They do not claim to exhaust theories outside those ledgers, such as modified
gravity, different visible-sector subtractions, interacting dark sectors,
non-perfect homogeneous residuals, anisotropic compact backreaction, non-FLRW
backgrounds, nonlocal effective stress tensors, or quantum Casimir sources.
\end{proof}


\section{Audit table and benchmark interface}

The preceding laws are intended to be operational.  A residual-closure audit can
be implemented without assuming any particular microscopic dark-sector model.  A
minimal background interface is
\begin{equation}
\begin{aligned}
&(E(a),\eps_{\vis}(a),p_{\vis}(a),\Omega_K)
\longrightarrow
(\eps_\D,p_\D)
\\
&\hspace{1.0cm}\longrightarrow
(\Aq,\Aw)
\longrightarrow
(\Rcl_q,\Rcl_{\wind},\mathfrak C_I).
\end{aligned}
\label{eq:standard_code_interface}
\end{equation}
This interface is deliberately less ambitious than a Boltzmann solver.  Its
purpose is to decide whether the adopted residual is even eligible to be called
minimal dust plus fixed vacuum before any topology, particle, or vacuum-origin
claim is made.

\begin{table*}[t]
\caption{Invariant residual-closure audit.  The table is a claim-control device:
each row must be satisfied at its own level before a stronger interpretation is
allowed.}
\label{tab:residual_audit_standard}
\begin{ruledtabular}
\begin{tabular}{p{0.24\textwidth}p{0.31\textwidth}p{0.34\textwidth}}
Audit question & Invariant object & Required statement \\
\hline
What is reconstructed? &
\(T^\D_{\mu\nu}=c^4G_{\mu\nu}/(8\pi G)-T^{\vis}_{\mu\nu}\) &
Metric, Einstein convention, and visible subtraction are declared. \\
Is the homogeneous residual dust plus vacuum? &
\((\Aq,\Aw)\), \(\mathfrak C_I\) &
\(\mathfrak C_I=0\) and \(\Aq,\Aw\ge0\) on the chosen interval. \\
Is dust microscopically realized? &
Brown current \(N^\mu\), energy law \(F(n)\) &
\(\nabla_\mu N^\mu=0\) and \(F(n)=\mu_{\wind}c^2n\). \\
Is the dust amplitude selected? &
\(\mu_{\wind}\), number rule, volume &
A non-topological mass scale and branch rule, such as \(N_J=|Q|\), are supplied. \\
Is vacuum equation of state realized? &
Four-form branch \(F_4\) &
Fixed branch with \(F_4=f\varepsilon_4\) and no unsuppressed membrane jump. \\
Is vacuum amplitude selected? &
\(\qhat\), \(\eF\), \(\NF\), branch prior &
Continuous value, discrete flux unit, or landscape/prior is specified. \\
Is topology predictive? &
\(\PG\), \(\Spec_\Gamma\), \(m_\Gamma(n)\), \({\rm Cov}_\Gamma\) &
Projector is propagated through transfer functions into observable covariance. \\
Is the likelihood comparison fair? &
\(\mathcal L_\Gamma-\mathcal L_{\rm bench}\) &
The benchmark has the same local background, amplitudes, mask, noise, and foreground conventions.
\end{tabular}
\end{ruledtabular}
\end{table*}

This audit table is the proposed standard.  It can be applied to a pure
\(\LCDM\) model, a \(w_0w_a\) dark-energy model, an interacting dark-sector
model, a compact-topology model, or a modified-gravity reconstruction.  The
standard does not decide which model is true; it prevents the model from claiming
more than its declared residual, microphysics, amplitude selector, and likelihood
actually establish.

A practical program follows three separable papers or software modules.  First,
construct the residual and closure-defect module Eq.~\eqref{eq:standard_code_interface}.  Second,
attach microscopic ledgers for current, vacuum, scalar-field, modified-gravity,
or interacting sectors.  Third, attach compact or noncompact transfer modules
that propagate initial statistics into observable covariance and likelihood.
Only the composition of all three modules is a predictive cosmological model.


\section{Conclusion}

The invariant content of the standard is the convention-declared residual
atlas
\begin{equation}
(g_{\mu\nu},T^{\vis}_{\mu\nu},\text{Einstein convention})
\longmapsto
T^\D_{\mu\nu}
\longmapsto
(\Aq,\Aw),
\end{equation}
with the affine subtraction law
\begin{align}
T^{\vis}_{\mu\nu}&\to T^{\vis}_{\mu\nu}+X_{\mu\nu},\\
(\Aq,\Aw)&\to
\bigl(\Aq+p_X,\,\Aw-a^3(\eps_X+p_X)\bigr).
\end{align}
At the homogeneous locally-FLRW level, this atlas is lossless.  The minimal
residual is dust-like matter plus fixed vacuum if and only if
\begin{equation}
-p_\D=\text{constant}\ge0,
\qquad
a^3(\eps_\D+p_\D)=\text{constant}\ge0 .
\end{equation}
Equivalently,
\begin{equation}
\eps_\D(a)=\eps_{\wind,0}a^{-3}+\eps_q,
\qquad
p_\D(a)=-\eps_q .
\end{equation}
This proves exact local background degeneracy with cold dark matter plus
\(\Lambda\), provided the same amplitudes, curvature, visible subtraction, and
perturbation covariance are used.

The natural minimal realization is a conserved Brown current with linear energy
law plus a fixed four-form vacuum branch:
\begin{equation}
F(n)=\mu_{\wind}c^2 n,
\qquad
T^q_{\mu\nu}=-\eps_q g_{\mu\nu}.
\end{equation}
The Brown current closes the dust equation of state, not the dust abundance.  The
four-form closes the vacuum equation of state, not the vacuum amplitude.  The
compact causal ledger closes null-reach bookkeeping on the closed cover or a
specified quotient, not the homogeneous stress tensor.

Thus the theory is invariantly admissible and background-closed at its stated
level, but not predictively complete until amplitude selection and compact
transfer likelihood are supplied.  The status ledger is:
\begin{description}
\item[Residual reconstruction:] closed.
\item[Homogeneous normal form:] closed.
\item[Dust equation of state:] closed by Brown current plus linear $F(n)$.
\item[Vacuum equation of state:] closed by fixed four-form branch.
\item[Dust amplitude:] not derived.
\item[Vacuum amplitude:] not derived.
\item[Local $R^3$ limit:] consistent.
\item[Compact causal ledger:] closed only after $A_0$, $\Gamma$, and $x_o$ are specified.
\item[Compact spectral closure:] requires quotient eigenmodes and projector.
\item[Perturbations:] source null tests stated; transfer closure incomplete.
\item[Observational likelihood:] not completed here.
\end{description}

The theorem also fixes the no-go boundary.  Topology may restrict the global
volume, causal reach, projected harmonic content, covariance matrix, and
observer-position prior.  It does not by itself provide the dimensionful current
mass scale, four-form flux unit, branch rule, or likelihood evidence required to
determine \(\Omega_c\), \(\Omega_\Lambda\), or \(\Lambda_{\eff}\).  Thus a
compact quotient is a background, causal, spectral, and covariance input, not an
abundance-origin mechanism.

The strongest form of the result is therefore a claim-discipline theorem:
every FLRW dark-sector residual that is exactly dust plus fixed vacuum is
equivalent, at the invariant background level, to a conserved Brown current with
linear energy law plus a constant four-form vacuum branch.  Compact causal
geometry supplies global causal and spectral closure ledgers, but not the
amplitudes.  A homogeneous equation-of-state departure is detected by
\(\Rcl_q\ne0\) or \(\Rcl_{\wind}\ne0\).  A source departure is detected by
nonzero rest-frame pressure, anisotropic stress, or sound speed.  A vacuum-branch
departure is detected by flux variation or membrane transitions.  An amplitude
claim requires non-topological normalization data.  A topology claim requires a
topology-covariance likelihood.  These alternatives exhaust the declared causal
residual compact-cosmology framework and make the manuscript self-contained at
its stated level.

\appendix

\section{Topology-likelihood diagnostics}
\label{app:transfer_diagnostics}

For each candidate $\Gamma$, generate mock skies from ${\rm Cov}^\Gamma$ and
matched controls from the benchmark covariance with the same local expansion
history, curvature, visible sector, amplitudes, mask, beam, foreground model,
and noise.  Hide the injection labels and run both pipelines.  With a threshold
$\Delta_\star$, report
\begin{align}
{\rm FPR}&=P(\Delta\ln\mathcal L>\Delta_\star|{\rm bench}),\\
{\rm FNR}&=P(\Delta\ln\mathcal L<\Delta_\star|\Gamma),\\
{\rm power}&=1-{\rm FNR} .
\end{align}
A topology claim is not likelihood-grade unless the pipeline distinguishes
quotient covariance from the benchmark covariance on blind synthetic data before
real-sky application.

For matched-circle and low-$\ell$ exclusions define
\begin{equation}
r_{\rm inj}(x_o)=\frac{1}{2}\min_{\gamma\ne e}d(x_o,\gamma x_o) .
\end{equation}
If $r_{\rm inj}(x_o)\ge\chi_*$, the last-scattering sphere predicts no matched
circles for that observer.  If $r_{\rm inj}(x_o)<\chi_*$, each holonomy with
$d(x_o,\gamma x_o)<2\chi_*$ predicts a circle pair.  In spherical geometry the
circle radius obeys
\begin{equation}
\cos\alpha_\gamma=
\frac{\tan[d(x_o,\gamma x_o)/2]}{\tan\chi_*},
\end{equation}
with relative phase/twist fixed by the holonomy.  Let
$S_{\rm circles}^\Gamma$ denote the resulting maximized circle statistic, and
let
\begin{equation}
\chi^2_{{\rm low}\text{-}\ell,\Gamma}
=d_{\rm low}^\dagger C_{{\rm low},\Gamma}^{-1}d_{\rm low}
\end{equation}
be the low-$\ell$ covariance statistic.  The admissible quotient set is
\begin{equation}
\begin{aligned}
\mathcal A_\Gamma=\{\,&\Gamma,x_o,R,\eta:
S_{\rm circles}^\Gamma<S_{\rm max},\\
&\chi^2_{{\rm low}\text{-}\ell,\Gamma}<\chi^2_{\rm max},
\quad
\mathcal L_\Gamma\ \hbox{acceptable}\,\} .
\end{aligned}
\label{eq:admissible_quotient_set_v16}
\end{equation}
If this set is empty, the chosen quotient implementation is excluded.  If it is
nonempty, the quotient is compact-compatible but not detected unless its
evidence relative to the matched benchmark remains significant after blind
injection calibration.



\begin{thebibliography}{99}

\bibitem{Wald}
R.~M. Wald, \emph{General Relativity} (University of Chicago Press, Chicago, 1984).

\bibitem{HawkingEllis}
S.~W. Hawking and G.~F.~R. Ellis, \emph{The Large Scale Structure of Space-Time}
(Cambridge University Press, Cambridge, 1973).

\bibitem{Wolf}
J.~A. Wolf, \emph{Spaces of Constant Curvature}, 6th ed. (AMS Chelsea,
Providence, 2011).

\bibitem{Hatcher}
A. Hatcher, \emph{Algebraic Topology} (Cambridge University Press, Cambridge, 2002).

\bibitem{BrownFluid}
J.~D. Brown, Action functionals for relativistic perfect fluids,
Class. Quantum Grav. \textbf{10}, 1579--1606 (1993).


\bibitem{Aurilia}
A. Aurilia, H. Nicolai, and P.~K. Townsend,
Hidden constants: the theta parameter of QCD and the cosmological constant of
\(N=8\) supergravity, Nucl. Phys. B \textbf{176}, 509 (1980).

\bibitem{HenneauxTeitelboim}
M. Henneaux and C. Teitelboim,
The cosmological constant as a canonical variable,
Phys. Lett. B \textbf{143}, 415 (1984).

\bibitem{BrownTeitelboim}
J.~D. Brown and C. Teitelboim,
Neutralization of the cosmological constant by membrane creation,
Nucl. Phys. B \textbf{297}, 787 (1988).

\bibitem{BoussoPolchinski}
R. Bousso and J. Polchinski,
Quantization of four-form fluxes and dynamical neutralization of the cosmological constant,
JHEP \textbf{06}, 006 (2000).

\bibitem{Bardeen}
J.~M. Bardeen, Gauge-invariant cosmological perturbations,
Phys. Rev. D \textbf{22}, 1882 (1980).

\bibitem{KodamaSasaki}
H. Kodama and M. Sasaki, Cosmological perturbation theory,
Prog. Theor. Phys. Suppl. \textbf{78}, 1 (1984).

\bibitem{Mukhanov}
V. Mukhanov, H. Feldman, and R. Brandenberger,
Theory of cosmological perturbations,
Phys. Rep. \textbf{215}, 203 (1992).

\bibitem{MaBertschinger}
C.-P. Ma and E. Bertschinger,
Cosmological perturbation theory in the synchronous and conformal Newtonian gauges,
Astrophys. J. \textbf{455}, 7 (1995).

\bibitem{Cornish}
N.~J. Cornish, D.~N. Spergel, G.~D. Starkman, and E. Komatsu,
Constraining the topology of the Universe,
Phys. Rev. Lett. \textbf{92}, 201302 (2004).

\bibitem{PlanckTopology}
P.~A.~R. Ade \emph{et al.} (Planck Collaboration),
Planck 2015 results. XVIII. Background geometry and topology,
Astron. Astrophys. \textbf{594}, A18 (2016), arXiv:1502.01593.

\bibitem{Planck2018}
N. Aghanim \emph{et al.} (Planck Collaboration),
Planck 2018 results. VI. Cosmological parameters,
Astron. Astrophys. \textbf{641}, A6 (2020), arXiv:1807.06209.

\bibitem{SahaLens}
S. Saha \emph{et al.} (COMPACT Collaboration),
Cosmic topology. Part Ic. Limits on lens spaces from circle searches,
JCAP \textbf{01}, 004 (2025), arXiv:2409.02226.


\bibitem{COMPACTPartI}
P. Petersen \emph{et al.} (COMPACT Collaboration),
Cosmic topology. Part I. Limits on orientable Euclidean manifolds from circle searches,
JCAP \textbf{01}, 030 (2023); Erratum: JCAP \textbf{04}, E01 (2024),
arXiv:2211.02603.

\bibitem{COMPACTIIa}
J.~R. Eskilt \emph{et al.} (COMPACT Collaboration),
Cosmic topology. Part IIa. Eigenmodes, correlation matrices, and detectability of orientable Euclidean manifolds,
JCAP \textbf{03}, 036 (2024), arXiv:2306.17112.

\bibitem{COMPACTIIb}
C.~J. Copi \emph{et al.} (COMPACT Collaboration),
Cosmic topology. Part IIb. Eigenmodes, correlation matrices, and detectability of non-orientable Euclidean manifolds,
arXiv:2510.05030.

\bibitem{COMPACTIIIb}
A. Samandar \emph{et al.} (COMPACT Collaboration),
Cosmic topology. Part IIIb. Eigenmodes and correlation matrices of spin-2 perturbations in orientable Euclidean manifolds,
JCAP \textbf{08}, 015 (2025), arXiv:2503.08671.

\bibitem{COMPACTIIc}
J. Noltmann \emph{et al.} (COMPACT Collaboration),
Cosmic topology. Part IIc. Detectability with non-standard primordial power spectrum,
arXiv:2602.15527.

\bibitem{NegroCasimir}
A. Negro \emph{et al.} (COMPACT Collaboration),
Quantum signatures of cosmic topology: How Casimir backreaction transmits isotropy violation,
arXiv:2603.12319.

\bibitem{DESIDR2}
M.~Abdul-Karim \emph{et al.} (DESI Collaboration),
DESI DR2 results II: Measurements of baryon acoustic oscillations and cosmological constraints,
Phys. Rev. D \textbf{112}, 083515 (2025), arXiv:2503.14738.

\end{thebibliography}

\end{document}
